\scanpage{41}
\mspage{\handwritten{9}}
% archive sheet number: 40
finie à groupe d'opérateurs $G_j$ (resp. $G'_j$). Posons, pour $i\handwritten{\in}I$, $i\handwritten{>}j$, $G_i=G_j\handwritten{\times}_{S_j}S_i$ et $G=G_j\handwritten{\times}_{S_j}S$, et définissons de même $G'_i$, $G'$, $H_i$,$H$, $H'_i$, et $H'$. Notons $\operatorname{Dihom}_{S\text{-gr.}}((G,\handwritten{\unsure{H|H'|\mathbb{H}|\mathbb{H}'}}),(\handwritten{G}',H'))$ l'ensemble des di-morphismes de $S$-groupes et de $S$-préschémas à groupe d'opérateurs du couple $(G,H)$ dans le couple $(G',H')$. Alors l'application canonique de $\underset{i\handwritten{\geqslant}j}{\underset{\handwritten{\longrightarrow}}{\lim}}{\unsure{\mathbb{D}|D}\mathrm{ihom}}_{S_i\text{-gr.}}((G_i,H_i),(G'_i,H'_i))$ dans $\operatorname{Dihom}_{S\text{-gr.}}((G,H),(G',H'))$ est bijective.

(ii) Dans la situation rappelée au début de (10.0), soient $G$ un $S$-groupe \deleted{\unsure{fini}} de présentation finie et $H$ un $S$-préschéma de présentation finie à groupe d'opérateurs $G$; il existe alors un indice $j\handwritten{\in}J$, un $S_j$-groupe de présentation finie $G_j$, un $S_j$-préschéma de présentation finie $H_j$ à groupe d'opérateurs $G_j$ et un di-isomorphisme de $S$-groupes et de $S$-préschémas à groupes d'opérateurs de $(G_j\handwritten{\times}_{S_j}S,H_j\handwritten{\times}_{S_j}S)$ sur$(G,H)$

\deleted{Proposition}\inserted{Définition} 10.15.- Etant donné un $S$-foncteur enngroupes $G$ et un $S$-foncteur $H$ à groupe d'opérate$\overset{\handwritten{\curvearrowleft}}{\underset{\handwritten{\curvearrowright}}{\mathrm{ru}}}$s $G$, on dit que $H$ est un \uline{$S$-foncteur en espaces homogènes sous $G$ pour la topologie $\handwritten{\mathbb{T}}$}, si le morphisme canonique $G\handwritten{\times}_S H\handwritten{\longrightarrow}H\handwritten{\times}_S H$ (défini par $(g,h)\handwritten{\rightsquigarrow}(g.h,h)$ pour $g\handwritten{\in}G(T)$, $h\handwritten{\in}H(T)$) est un épimorphisme de faisceaux pour la topologie $\handwritten{\mathbb{T}}$.

Si $G$ est un $S$-groupe et si $H$ est un $S$-préschéma à \deleted{groupe} groupe d'opérateurs $G$, dire que $H$ est un $S$-préschéma en espaces homogènes sous $G$ pour la topologie $\handwritten{\mathbb{T}}$ revient donc à dire que le morphisme canonique $G\handwritten{\times}_S H\handwritten{\longrightarrow}H\handwritten{\times}_S H$ défini comme précédemment est couvrant pour la topologie $\handwritten{\mathbb{T}}$ (\deleted{Exp.} IV 4.4.3).

De même, dire qu'un $S$-préschéma $H$ à groupe d'opérateurs un $S$-groupe $G$ est un fibré principal homogène sous $G$ pour la topologie $\handwritten{\mathbb{T}}$ (\deleted{Exp.} IV 5.1.5) revient à dire (\deleted{Exp.} IV 5.6 (ii)) que le morphisme canonique $G\handwritten{\times}_S H\handwritten{\longrightarrow}H\handwritten{\times}_S H$ est un isomorphisme et que le morphisme structural $H\handwritten{\longrightarrow}S$ est couvrant pour la topologie $\handwritten{\mathbb{T}}$.

\uline{10.16.- Dans la situation rappelée au début de (10.0) , soient $j\handwritten{\in}I$, $G_j$ un $S_j$-grou-}
% The final word continues on the next scan in this package.
